\section{Plano de inteligencia de amenazas}
\label{sec:cti-impl}

El plano de inteligencia cumple dos funciones. La primera es dar a todo el sistema una
\textbf{referencia común}: las reglas, los patrones del Pilar~A, el grafo del Pilar~B y la
extracción de técnicas hablan en identificadores de ATT\&CK, y esos identificadores tienen que
significar lo mismo en todas partes. La segunda es convertir en datos estructurados la
inteligencia que se publica como texto, mediante la extracción automática de técnicas (ML-2).

\subsection{ATT\&CK como referencia común}

Se fija \textbf{ATT\&CK Enterprise v19.2} \cite{mitre_attack}, con su suma SHA-256 registrada. La
versión importa: la v19 dividió la táctica de evasión de defensas en dos y revocó técnicas, de
modo que un mismo identificador puede no significar lo mismo en dos versiones. La v20, publicada
durante el desarrollo y parcheada el mismo día, se descartó; subir de versión es un paso aparte
que se medirá.

La carga en el grafo usa un \textbf{cargador STIX propio e idempotente}: la identidad de cada nodo
es su identificador STIX, cada nodo y relación guarda la lista de cargas que lo aportaron, y las
fases de la \textit{kill chain} de cada técnica se convierten en aristas hacia su táctica. El mismo
cargador sirve para ATT\&CK, para un feed de indicadores y para los perfiles que genera ML-2.
ATT\&CK queda en el grafo con 4.823 nodos y 22.352 relaciones en 4,7\,s, y recargarlo no
duplica nada.

Sobre el mismo paquete se construye un \textbf{servicio de normalización de técnicas}. Dado un
identificador, sigue las cadenas de revocación hasta el vigente (por ejemplo,
\texttt{T1086}~$\rightarrow$~\texttt{T1059.001}), marca los obsoletos sin sustituto y sube de
subtécnica a técnica padre. Medido contra la v19.2, las alertas y campañas del Pilar~A están
vigentes al 100\,\%; en la verdad conocida hay dos técnicas revocadas, ninguna de ellas usada por
los patrones.

Como \textbf{feed de indicadores} se incorpora ThreatFox \cite{threatfox}, convertido de CSV a STIX
con identificadores deterministas: 135.737 indicadores de 508 familias de malware, cargados en
17\,s. Cada familia se enlaza con el software de ATT\&CK por nombre o alias \textbf{solo si la
coincidencia es única}; las ambiguas se descartan (por ejemplo, Snake también coincidía con el
Uroburos de Turla). Así, el 25,6\,\% de los indicadores llega a técnicas de ATT\&CK a través de su
familia. El resto pertenece a familias que no están en el catálogo de software de ATT\&CK: el
límite es el catálogo, no el cruce.

\subsection{Plataforma CTI frente a carga directa}

Se comparó OpenCTI \cite{opencti}, la plataforma de referencia citada en el
Capítulo~\ref{cap:estadoarte}, con la carga directa descrita, sobre los mismos datos y con
criterios fijados antes de medir (Tabla~\ref{tab:opencti}).

\begin{table}[H]
\centering
\begin{tabular}{|l|r|r|}
\hline
\textbf{Criterio} & \textbf{Carga directa} & \textbf{OpenCTI} \\
\hline
Ingesta de ATT\&CK (22.600 objetos) & 4,7\,s & 32\,min \\
Ingesta de ThreatFox (272.853 objetos) & 17\,s & 3\,h 28\,min \\
Memoria con ATT\&CK y ThreatFox & 3,6\,GB & 13,7\,GB \\
Disco con ATT\&CK y ThreatFox & 865\,MB & 1,9\,GB \\
Servicios & 1 & 11 \\
\hline
\end{tabular}
\caption{OpenCTI frente a la carga directa en el grafo.}
\label{tab:opencti}
\end{table}

El coste no es lo más relevante; la \textbf{fidelidad}, sí. OpenCTI descarta las 149 técnicas
revocadas y sus relaciones de sustitución, que son justo lo que permite traducir identificadores
antiguos. Además fusiona grupos y malware que comparten un alias (por ejemplo, Ember Bear con Saint
Bear) y no importa las estrategias de detección ni las analíticas de ATT\&CK v18 en adelante.
Conserva intactas las relaciones de uso, mitigación y atribución. Por eso \textbf{la referencia
del sistema es la carga directa}, y OpenCTI queda como capa opcional para el analista y para la
gestión de indicadores, cuyo valor añadido se mide aparte.

\subsection{Extracción de técnicas de informes (ML-2)}
\label{sec:ml2}

\medskip\noindent\textbf{Tarea y datos.} Se usa AnnoCTR \cite{annoctr}, con 120 informes de CTI anotados por
expertos. Cada mención de una técnica está enlazada con ATT\&CK, y dos tercios de las menciones son
\textbf{implícitas}: describen la técnica sin nombrarla. La tarea se plantea como clasificación
multietiqueta \textbf{frase~$\rightarrow$~técnicas padre}, con las particiones publicadas, que
separan por informe (70, 16 y 34 informes para entrenamiento, desarrollo y prueba). El espacio
de etiquetas lo forman las 59 técnicas con al menos cinco frases de entrenamiento, que cubren el
85\,\% de las menciones de prueba. Las frases sin técnica se mantienen como ejemplos negativos.
Se evalúa a tres niveles: por frase, por mención (explícita o implícita) y por informe, que es
el nivel del perfil que se carga en el grafo.

\medskip\noindent\textbf{Sistemas.}
\begin{itemize}
    \item \textbf{Línea base de diccionario:} empareja los nombres e identificadores de ATT\&CK. Los
    nueve términos genéricos que más fallan (como «Malware» o «DNS») se podan usando solo la
    partición de desarrollo.
    \item \textbf{SecureBERT~2.0} \cite{securebert2}, preentrenado en textos de ciberseguridad, frente
    a \textbf{ModernBERT-base} \cite{modernbert}, el modelo general del que parte. Como comparten
    arquitectura, la diferencia entre ambos aísla el efecto del preentrenamiento en el dominio. El
    SecureBERT original solo publica sus pesos en un formato que puede ejecutar código al
    cargarse, por lo que se descartó. Solo se cargan pesos en \texttt{safetensors} y con revisión
    fijada.
\end{itemize}

\medskip\noindent\textbf{Protocolo.} La época y el umbral de decisión se eligen en desarrollo, y la partición
de prueba se evalúa una sola vez. Cada modelo se entrena con tres semillas, cada entrenamiento
tarda unos cinco minutos en la GPU de la estación y todo queda registrado en MLflow. Las
diferencias entre sistemas se contrastan con un \textbf{remuestreo emparejado por informe}
\cite{efron1993,bergkirkpatrick2012}, con 10.000 muestras. Se remuestrea por informe porque las
frases de un mismo informe no son independientes.

\medskip\noindent\textbf{Resultados.} La Tabla~\ref{tab:ml2} resume la partición de prueba. Ambos modelos superan
con claridad al diccionario, y SecureBERT~2.0 supera a ModernBERT de forma significativa:
+0,034 de F1 por frase, con un intervalo de confianza del 95\,\% de [+0,011, +0,058]. El
preentrenamiento en el dominio aporta. La ventaja se concentra en las menciones explícitas; en
las implícitas, los dos modelos casi empatan.

\begin{table}[H]
\centering
\resizebox{\textwidth}{!}{%
\begin{tabular}{|l|c|c|c|c|c|}
\hline
\textbf{Sistema (prueba)} & \textbf{F1 frase} & \textbf{F1 macro} & \textbf{R explícitas} & \textbf{R implícitas} & \textbf{F1 informe} \\
\hline
Diccionario & 0,330 & 0,349 & 0,355 & 0,151 & 0,490 \\
ModernBERT-base & 0,519 & 0,262 & 0,670 & 0,411 & 0,620 \\
SecureBERT 2.0 & 0,553 & 0,301 & 0,721 & 0,424 & 0,670 \\
\textbf{SecureBERT 2.0 $\cup$ diccionario} & \textbf{0,590} & \textbf{0,464} & \textbf{0,832} & \textbf{0,488} & \textbf{0,713} \\
\hline
\end{tabular}%
}
\caption{ML-2 en la partición de prueba de AnnoCTR (media de tres semillas; R: exhaustividad
por mención).}
\label{tab:ml2}
\end{table}

El análisis por técnica explica por qué conviene combinar modelo y diccionario. El F1 de cada
técnica crece con su frecuencia en entrenamiento (correlación de Spearman de 0,61): el modelo gana
en las técnicas con veinte o más frases de entrenamiento, y el diccionario en las raras cuando
aparecen nombradas. Elegida en desarrollo, sin embargo, la unión de los dos perdía 0,011. La causa
es que la partición de desarrollo apenas nombra técnicas: el diccionario rinde allí la tercera
parte que en prueba. Con solo 16 informes no daba para decidir, y \textbf{la decisión se tomó con
una validación cruzada por informe sobre los 86 de entrenamiento y desarrollo}, con reglas fijadas
antes de ejecutarla. La unión mejora en +0,040 de F1 por frase, con un intervalo de [+0,020,
+0,060], y gana en cuatro de los cinco bloques. Es el sistema adoptado. La lección vale para
cualquier conjunto pequeño: con 16 informes se habría concluido justo lo contrario.

La precisión medida está además \textbf{infraestimada}. En una revisión manual de 60 de los 270
falsos positivos (Anexo~\ref{anx:ml2}), el 42\,\% resultaron técnicas que la frase sí describe y
que AnnoCTR no anotó, como filas de tablas de ATT\&CK u órdenes de PowerShell literales. Contándolas
como aciertos, la precisión por frase pasaría de 0,61 a unos 0,77. Los errores reales se deben a
técnicas vecinas de una misma cadena (intérprete de órdenes frente a \textit{phishing}, ofuscación
frente a canal cifrado) y a tomar por comportamiento del adversario la discusión general de una
tecnología.

\medskip\noindent\textbf{Evaluación sobre el corpus del proyecto.} El sistema se aplicó a seis informes públicos
de Microsoft sobre APT29, de los que se descartan las secciones de defensa y de producto. El perfil
resultante se compara con la ficha de ATT\&CK del grupo y con lo que detectó el propio sistema.

\begin{itemize}
    \item \textbf{Frente a la ficha de ATT\&CK} (G0016 y sus dos campañas): de sus 72 técnicas,
    solo 33 están en el espacio de etiquetas. Frente a esas 33, el perfil alcanza un F1 de 0,69,
    con precisión 0,83 y exhaustividad 0,60. El diccionario solo llega a 0,49 y el modelo solo a
    0,53.
    \item \textbf{Frente a lo detectado en el primer día de APT29}: dentro del espacio de etiquetas,
    el perfil extraído cubre 4 de las 5 técnicas de las campañas detectadas, igual que la ficha
    oficial del grupo. La laguna más relevante es estructural: \textbf{el volcado de credenciales
    (\texttt{T1003}), núcleo del patrón P1, no está entre las 59 etiquetas}, porque el espacio sale
    de la frecuencia en AnnoCTR y no de lo que detecta el sistema.
    \item \textbf{Parecido entre campaña y actor}: se ordenaron los 173 grupos de ATT\&CK por parecido
    de técnicas con las campañas detectadas. G0016 no aparece arriba con ninguna de las medidas
    fijadas: queda entre el puesto 39 y el 85, y ni siquiera con la verdad del plan de emulación
    como consulta pasa del 57. Las campañas tienen pocas técnicas, genéricas y compartidas por
    muchos grupos. Es un resultado negativo que se informa como tal: \textbf{comparar conjuntos de
    técnicas no basta para reconocer a un actor}.
\end{itemize}

Este último bloque tiene una circularidad que se declara: la ficha de G0016 se construye en parte
con estos mismos informes, y AnnoCTR incluye informes de APT29, de otros editores, en
entrenamiento. Coincidir con la ficha es esperable y no prueba que el modelo generalice; eso lo
prueba la evaluación sobre AnnoCTR. El perfil se exporta en STIX~2.1, con la confianza, la
evidencia y el método de cada técnica, y se carga en el grafo con el mismo cargador que ATT\&CK.
